%! AP Statistics — Unit 4, Lesson 4.12 — Student Packet Cover Sheet
\documentclass[10pt]{article}
\usepackage{apstats-article}
\usepackage{apstats-boxes}
\usepackage{ltablex}
\keepXColumns

\begin{document}

% Full-bleed navy banner
\begin{tikzpicture}[remember picture, overlay]
  \fill[navy] (current page.north west)
    rectangle ([yshift=-0.9in]current page.north east);
\end{tikzpicture}
\vspace{-0.8in}

\begin{center}
  {\color{white}\textbf{\LARGE AP Statistics}}\\[4pt]
  {\color{skymid}\large\textbf{Unit 4: Probability, Random Variables, and Probability Distributions}}\\[6pt]
  {\color{white}\textbf{\large Lesson 4.12 \quad The Geometric Distribution}}
\end{center}

\vspace{0.22in}
\namedateperiod
\vspace{0.18in}

\begin{learningtargetbox}
\small
\begin{enumerate}[leftmargin=1.5em, itemsep=5pt]
  \item \ldots{} determine whether a scenario is a geometric setting by checking
        the BITS conditions, and compute $P(X = k) = (1-p)^{k-1}p$ using the
        formula and \texttt{geometpdf}.
  \item \ldots{} find and interpret the mean $\mu_X = 1/p$ of a geometric
        distribution, and compute cumulative probabilities with \texttt{geometcdf}.
\end{enumerate}
\end{learningtargetbox}

\vspace{0.14in}

\begin{tocbox}
\small
\renewcommand{\arraystretch}{1.5}
\begin{tabularx}{\linewidth}{@{} c l X r @{}}
\rowcolor{sky}
\textbf{\#} & \textbf{Component} & \textbf{Description} & \textbf{Score} \\
\midrule
1 & Warm-Up        & Spiral review: binomial parameters; contrast with geometric  & \blank{1.2cm} \\
2 & Guided Notes   & BITS; geometric formula; $\mu_X = 1/p$; \texttt{geometpdf/cdf} & \blank{1.2cm} \\
3 & Group Activity & Geometric vs.\ binomial; computing probabilities             & \blank{1.2cm} \\
4 & Exit Ticket    & Assembly-line defect inspection                              & \blank{1.2cm} \\
5 & Homework       & Four geometric scenarios; AP free-response                   & \blank{1.2cm} \\
\midrule
  & \multicolumn{2}{r}{\textbf{Total}} & \blank{1.2cm} \\
\end{tabularx}
\end{tocbox}

\vspace{0.14in}

\begin{remindbox}
\small
The \textbf{most important contrast in Lesson 4.12}: geometric and binomial
both require binary, independent, same-$p$ trials --- but they count different things.

\vspace{6pt}
\begin{tabularx}{\linewidth}{@{} >{\centering\arraybackslash\columncolor{goldbg}}p{0.13\linewidth}
                                    X @{}}
\textbf{Binomial} &
$n$ is fixed in advance. $X$ counts the number of successes in $n$ trials.
Formula: $P(X=k) = \binom{n}{k}p^k(1-p)^{n-k}$. Mean: $\mu_X = np$. \\[6pt]
\textbf{Geometric} &
$n$ is NOT fixed. $X$ counts the number of trials up to and \emph{including}
the first success. Formula: $P(X=k) = (1-p)^{k-1}p$. Mean: $\mu_X = 1/p$. \\[6pt]
\textbf{Key:} &
In the geometric formula, the exponent is $k-1$ (the $k-1$ failures before
the success on trial $k$), not $k$.
\end{tabularx}
\end{remindbox}

\end{document}
